% Requires amsmath, amssymb, hyperref.
\section{The odd-prime estimate for the minus part}
\label{sec:odd-prime-minus}

The arithmetic window of \eqref{eq:entry-strict-length} bounds the full
strict cohomology group. One of the estimates contributing to that
bound is an inequality for its minus part at all odd primes. Two
different calculations enter the inequality. Character valuations
control the primes away from the degree of the auxiliary field. An
annihilator controls the primes dividing the degree. The sums below
retain their factors of \(\log\ell\) throughout.

\subsection{Statement for a fixed field and range of twists}

Write \(E\) for the auxiliary CM field, \(D=[E:\mathbb Q]\), and
\(F\) for its conductor, the least positive integer such that \(E\)
embeds in \(\mathbb Q(\zeta_F)\). For an odd prime \(\ell\) and
\(k\ge1\), let \(m_{\ell,k}\) be the \(\mathbb Z_\ell\)-length of the
minus part of the strict group for the twist \(-k\). More precisely,
if \(c\) acts on the twisted cohomology by complex conjugation, the
minus part is \(T_{\ell,k}(E)\cap\ker(c-(-1)^{k+1})\). The plus part
has eigenvalue \((-1)^k\). Conjugation preserves the strict kernel and
squares to the identity. Since \(2\) is a unit in \(\mathbb Z_\ell\),
the two operators \((1+(-1)^kc)/2\) and \((1-(-1)^kc)/2\) restrict to
complementary idempotents of the strict kernel. They give a direct sum
of these subgroups and hence additivity of the two lengths, even
before finiteness has been established.
\footnote{\href{https://github.com/Deicyde/riemann-hypothesis/blob/40e56bab/src/Development/AutomorphicArithmetic/AutomorphicVarieties/LengthEqAddEigen.lean\#L68-L103}{Parity and length additivity for odd primes}.}

Fix a positive integer \(m\), a subgroup
\(H\subseteq(\mathbb Z/m\mathbb Z)^\times\), and put
\(d_0=[(\mathbb Z/m\mathbb Z)^\times:H]\). Choose \(e\ge2\) and
primes \(q\equiv3\bmod8\), \(d\equiv7\bmod16\) with
\(\gcd(qd,m)=1\). In the field construction of
Section~\ref{sec:entry-assembly}, replace the test \(\psi(u)=1\)
by \(u\bmod m\in H\). Then
\[
       \ell\mid D,\quad \ell\ne2\quad\Longrightarrow\quad
       \ell\mid d_0
       \qquad(\ell\text{ prime}).
\]
To prove this divisibility, choose the cyclotomic field in the
definition of \(E\), and let \(M\) be the order of its chosen
primitive root. Each automorphism raises the root to a power \(u\)
prime to \(M\). Let \(h=2^{e+1}d_0\). Since
\(\varphi(2^{e+2})=2^{e+1}\), we have \(u^h\equiv1\bmod 2^{e+2}\).
The class of \(u\) modulo \(H\) has order dividing \(d_0\), so
\(u^h\bmod m\in H\). The exponent \(h\) is even and hence both
quadratic symbols of \(u^h\) are \(1\). It follows from the
fixed-field definition that the \(h\)-th power of every cyclotomic
automorphism fixes \(E\). Its fixing subgroup is normal since the
cyclotomic Galois group is abelian, and the quotient has order
\([E:\mathbb Q]=D\). For any prime \(\ell\mid D\), Cauchy's theorem
now gives a quotient element of order \(\ell\). Its \(h\)-th power is
the identity, so \(\ell\mid h\). When \(\ell\) is odd this implies
\(\ell\mid d_0\), as asserted.
\footnote{\href{https://github.com/Deicyde/riemann-hypothesis/blob/40e56bab/src/Development/AutomorphicArithmetic/AutomorphicVarieties/StmtFields.lean\#L11-L45}{The cyclotomic field and subgroup data},
\href{https://github.com/Deicyde/riemann-hypothesis/blob/40e56bab/src/Development/AutomorphicArithmetic/AutomorphicVarieties/OddResidueBounds/DegExistsLift.lean\#L33-L99}{Powering in the cyclotomic quotient and the degree divisibility}.}

There is an absolute \(C\ge0\) from the annihilator estimate
such that, for every integer \(t\ge2\), every finite set \(S\) of
odd primes, and every allowed choice of local data, the lengths
\(m_{\ell,k}\) with \(1\le k<t\) are finite and
\begin{align}
 \sum_{\ell\in S}(\log\ell)\sum_{k=1}^{t-1}m_{\ell,k}
 &\le 3Dt^2\log(tF) \notag\\
 &\quad{}+
 C\sum_{\substack{\ell\mid d_0\\\ell\ {\rm odd\ prime}}}
       \ell D^2t^2\bigl(\log(t\ell F)\bigr)^2 .
 \label{eq:galois-odd-minus-bound}
\end{align}
The same \(C\) works for all the fields and all the finite sets of primes.
\footnote{\href{https://github.com/Deicyde/riemann-hypothesis/blob/40e56bab/src/Development/AutomorphicArithmetic/AutomorphicVarieties/OddResidueBounds/ThmBcOfPsiBound.lean\#L147-L177}{The field conductor and the split according to its degree},
\href{https://github.com/Deicyde/riemann-hypothesis/blob/40e56bab/src/Development/AutomorphicArithmetic/AutomorphicVarieties/OddResidueBounds/ThmBcOfPsiBound.lean\#L266-L295}{the uniform odd-prime conclusion and its inputs}.}

\subsection{Primes dividing the degree}

For an odd \(\ell\) the annihilator input gives a nonnegative integer
\(r\) killing the minus subgroup by multiplication by \(\ell^r\), with
\[
       r\log\ell\le C_0D(k+1)\log(k\ell F)
\]
for one absolute \(C_0\ge0\). A second input gives an integer \(g\)
bounding the length of the subgroup killed by \(\ell\) in the
ambient global \(H^1\), with \(g\le 8\ell D\log(\ell F)\).

To see exactly how these estimates combine, let \(P\) be the minus
subgroup and \(P_i=P[\ell^i]\). Its filtration starts with \(P_0=0\)
and ends with \(P_r=P\). Multiplication by \(\ell^{i-1}\) embeds
\(P_i/P_{i-1}\) into the \(\ell\)-torsion of the ambient global group.
Indeed its image is killed by \(\ell\), and the kernel on \(P_i\) is
exactly \(P_{i-1}\). Thus \(\operatorname{length}P\le rg\), including
the case \(r=0\). This proves finiteness. Since all factors are
nonnegative and \(\log(\ell F)\le\log(k\ell F)\), it also gives
\begin{align}
 m_{\ell,k}\log\ell
 &\le (r\log\ell)g \notag\\
 &\le 8C_0\ell D^2(k+1)\bigl(\log(k\ell F)\bigr)^2 .
 \label{eq:galois-annihilator-length}
\end{align}
Set \(C=8C_0\). This estimate actually applies at every odd \(\ell\).
Using it only for \(\ell\mid D\) keeps its factor of \(\ell\) out
of the unbounded prime sum.\footnote{\href{https://github.com/Deicyde/riemann-hypothesis/blob/40e56bab/src/Development/AutomorphicArithmetic/AutomorphicVarieties/OddResidueBounds/ThmBcOfPsiBound.lean\#L15-L39}{The annihilator, generator estimate and resulting finite length with constant \(8C_0\)}.}

\subsection{The character sum away from the degree}

Let \(\rho:(\mathbb Z/F\mathbb Z)^\times\to
\operatorname{Gal}(E/\mathbb Q)\) be the cyclotomic action. It is
surjective, and \([\left(\mathbb Z/F\mathbb Z\right)^\times:\ker\rho]=D\).
For a fixed \(k\), let \(\mathcal C\) be the set of complex characters
trivial on \(\ker\rho\) such that \(\chi(-1)=(-1)^{k+1}\).
Passing to the quotient embeds
\(\mathcal C\) in the character group of
\((\mathbb Z/F\mathbb Z)^\times/\ker\rho\). The complex character
group of a finite abelian group has its order, so
\(\lvert\mathcal C\rvert\le D\).

Here are the valuation estimates used for this step. For
\(\ell\nmid 2D\) the characterwise input gives finiteness of
\(m_{\ell,k}\) and a finite set \(T_\ell\) of allowed \(\ell\)-adic
characters with
\[
     m_{\ell,k}\log\ell
       \le\sum_{\chi\in T_\ell}-\log
              \lvert\alpha_F(k+1,\chi)\rvert_\ell.
\]
To specify the normalization, write \(f\) for the conductor of \(\chi\)
and \(\chi^*\) for the primitive character inducing it. With \(B_j(X)\)
the \(j\)-th Bernoulli polynomial, the scalar is
\[
    \alpha_F(j,\chi)=F\!\!\prod_{\substack{p\ {\rm prime}\\p\le j+1}}\!p
          \,\cdot f^{j-1}\sum_{a=1}^f\chi^*(a)\,B_j(a/f).
\]
The product-valuation estimate gives nonnegative integers \(w_\ell\),
the same ones for
all these choices, such that
\begin{align*}
  \sum_{\chi\in T_\ell}-\log
              \lvert\alpha_F(k+1,\chi)\rvert_\ell
       &\le w_\ell\log\ell,\\
  \sum_{\ell\in S'}w_\ell\log\ell
       &\le\lvert\mathcal C\rvert\log B(F,k+1),&
 B(F,j)&=16F^{j+1}j!
\end{align*}
for each finite prime set \(S'\). Applying the first inequality to
each member of an arbitrary \(S'\) disjoint from the divisors of
\(2D\), and then using the second, yields
\begin{equation}\label{eq:galois-good-prime-sum}
      \sum_{\ell\in S'}m_{\ell,k}\log\ell
          \le D\log B(F,k+1).
\end{equation}
The last inequality uses \(B(F,k+1)\ge1\). The characterwise input and
the product-valuation estimate are the two arithmetic premises of
this deduction.\footnote{\href{https://github.com/Deicyde/riemann-hypothesis/blob/40e56bab/src/Development/AutomorphicArithmetic/AutomorphicVarieties/OddResidueBounds/ThmBcOfPsiBound.lean\#L41-L108}{The quotient-character injection, valuation estimates and sum over primes},
\href{https://github.com/Deicyde/riemann-hypothesis/blob/40e56bab/src/Development/AutomorphicArithmetic/AutomorphicVarieties/OddResidueBounds/ZeroMemUnrHom.lean\#L79-L110}{the characterwise arithmetic estimates},
\href{https://github.com/Deicyde/riemann-hypothesis/blob/40e56bab/src/Development/AutomorphicArithmetic/AutomorphicVarieties/GaloisLibrary/BernoulliNumbers/MapGenBChar.lean\#L9-L31}{Bernoulli and denominator conventions}.}

\subsection{Summing over the twists}

If \(1\le k<t\), then \(tF\ge2\) and \(k+1\le t\). Consequently
\[
  16\le(tF)^4,\qquad F^{k+2}\le(tF)^{t+1},\qquad
  (k+1)!\le(k+1)^{k+1}\le(tF)^t.
\]
Multiplying and taking logarithms gives
\(\log B(F,k+1)\le(2t+5)\log(tF)\). For the primes in \(S\) away
from \(D\), interchange the two finite sums in
\eqref{eq:galois-odd-minus-bound} and apply
\eqref{eq:galois-good-prime-sum} at each \(k\). The result is at most
\[
     D(t-1)(2t+5)\log(tF)\le 3Dt^2\log(tF),
\]
since \(3t^2-(t-1)(2t+5)=t^2-3t+5>0\).

For a remaining prime use \(k+1\le t\) and
\(0\le\log(k\ell F)\le\log(t\ell F)\) in
\eqref{eq:galois-annihilator-length}. Its contribution for all twists is
at most \(C\ell D^2(t-1)t(\log(t\ell F))^2\), hence at most
\(C\ell D^2t^2(\log(t\ell F))^2\). Every such \(\ell\) is an odd
prime divisor of \(d_0\). Enlarging the sum to all these divisors
only adds nonnegative terms. Adding the two bounds proves
\eqref{eq:galois-odd-minus-bound}. The proof applies at each fixed
field and range of twists, so it can be applied at \(e(x),q,d\) and
\(t=\lfloor x^{401}\rfloor\) as soon as those data meet the field
construction's hypotheses.\footnote{\href{https://github.com/Deicyde/riemann-hypothesis/blob/40e56bab/src/Development/AutomorphicArithmetic/AutomorphicVarieties/OddResidueBounds/ThmBcOfPsiBound.lean\#L110-L145}{The factorial and conductor estimate},
\href{https://github.com/Deicyde/riemann-hypothesis/blob/40e56bab/src/Development/AutomorphicArithmetic/AutomorphicVarieties/OddResidueBounds/ThmBcOfPsiBound.lean\#L178-L289}{summing the two ranges}.}

In the passage to \eqref{eq:entry-strict-length} this odd-prime
estimate is combined with the dyadic and fixed-prime estimates, the
other cohomological counts and the global window construction. The
source applies the odd-prime theorem to its characterwise comparison
obtained from Poitou--Tate duality and the supplied odd-prime
input.\footnote{\href{https://github.com/Deicyde/riemann-hypothesis/blob/40e56bab/src/Development/AutomorphicArithmetic/AutomorphicVarieties/ModularityArgument/ThresholdArguments/DenseThresholdCases/HVCore.lean\#L18-L22}{Application inside the proof of the full window},
\href{https://github.com/Deicyde/riemann-hypothesis/blob/40e56bab/src/Development/AutomorphicArithmetic/AutomorphicVarieties/ModularityArgument/ThresholdArguments/DenseThresholdCases/WindowBoundOfLiteratureLevel.lean\#L25-L34}{the automorphic hypotheses furnishing the other counts}.}

\section{The conjugate weights of the global Galois representation}
\label{sec:galois-parallel-weight}

Let \(E/E^+\) be a CM field with maximal totally real subfield \(E^+\)
and complex conjugation \(c\). Fix \(a\ge1\) and let \(\pi\) be regular
algebraic and cuspidal on \(\mathrm{GL}_a\) over \(E\), polarized by a
Hecke character \(\chi\) of \(E^+\). For a complex embedding \(\sigma\)
let \(I_\sigma\) be the multiset of infinitesimal-character parameters
of \(\pi\). We derive the integer \(b\) and the multisets \(H_\tau\)
required by Section~\ref{sec:entry-assembly}. The relation with the
actual jumps of the Galois representation uses the local de Rham
comparison at the end of the argument.

\subsection{The exponent at infinity of the polarizing character}

Let \(\psi=\chi\circ\operatorname{N}_{E/E^+}\), and write \(s_\sigma\)
for the exponent of its archimedean character at \(\sigma\). The
infinitesimal-character rules for twisting and contragredients, applied
to the invariant pairing expressing the polarization, give
\begin{equation}\label{eq:galois-inf-char}
    I_{\sigma c}=\{-(z-s_\sigma):z\in I_\sigma\}.
\end{equation}
Here twisting by \(\psi^{-1}\) adds \(-s_\sigma\), and passage across
the pairing negates the resulting parameter. The multiset \(I_\sigma\)
has cardinality \(a\ge1\). Take one member \(z\). Regular algebraicity
gives \(z=(a-1)/2+n\), and likewise
\(s_\sigma-z=(a-1)/2+n'\), for integers \(n,n'\). Adding shows
\[
                  s_\sigma=a-1+n+n'\in\mathbb Z.
\]
The pullback by the norm gives \(s_{\sigma c}=s_\sigma\). Hence these
integers give a well-defined \(s_\phi\) indexed by the real embeddings
\(\phi:E^+\to\mathbb R\).

For \(x\in(E^+)^\times\) totally positive, put
\(u_\phi=\exp(\tfrac12\log\phi(x))\) and regard \(u\) as a positive
idele at the complex places of \(E\), equal to \(1\) at finite places.
Its norm to \(E^+\) has archimedean components \(\phi(x)\). Evaluating
the archimedean character of \(\psi\) at \(u\) multiplies the two
equal contributions above each real embedding. It follows that the
archimedean component of \(\chi\) satisfies
\begin{equation}\label{eq:galois-inf-weight}
              \chi_\infty(x)=\prod_\phi\phi(x)^{s_\phi}.
\end{equation}
\footnote{\href{https://github.com/Deicyde/riemann-hypothesis/blob/40e56bab/src/Development/AbelianVarietyTheory/Proofs/SatakeParameters/ExistsIsWeight.lean\#L25-L67}{The invariant pairing, integrality and conjugation of the exponents},
\href{https://github.com/Deicyde/riemann-hypothesis/blob/40e56bab/src/Development/AbelianVarietyTheory/Proofs/SatakeParameters/ExistsIsWeight.lean\#L67-L138}{the norm at the archimedean places}.}

In fact \(s_\phi\) is independent of \(\phi\). To see this in the form
used in the construction, take a unit \(v\) of \(E^+\). The finite
idele of \(v\), raised to the powers \(j!\), converges to \(1\).
Indeed the image of \(v\) in the unit group of every finite quotient
by a nonzero ideal has finite order, which divides \(j!\) eventually.
These congruences prove convergence in the product of the local unit
groups. Continuity of \(\chi\) shows that the modulus of its value
on this finite idele is \(1\), since a positive real \(r\) with
\(r^{j!}\to1\) must be \(1\). The value on the full principal idele
is \(1\), so \(\lvert\chi_\infty(v)\rvert=1\). Apply
\eqref{eq:galois-inf-weight} to the totally positive element \(v^2\)
and take logarithms to obtain
\[
               \sum_\phi s_\phi\log|\phi(v)|=0.
\]
By the product formula also \(\sum_\phi\log|\phi(v)|=0\). Choose
\(\phi_0\) minimizing \(s_\phi\). In degree one the exponents are
already constant. In higher degree, Dirichlet's unit theorem gives a
unit \(v\) with \(\log|\phi(v)|<0\) for all \(\phi\ne\phi_0\), since
the full log-unit lattice in the sum-zero hyperplane meets the open
cone given by those strict inequalities. Subtracting \(s_{\phi_0}\) times the
product formula from the preceding equality for this \(v\) gives
\[
         \sum_\phi (s_\phi-s_{\phi_0})\log|\phi(v)|=0 .
\]
Every summand is nonpositive, and a strict increase in any of the
integers \(s_\phi\) would make its summand negative. Thus there is
one integer \(s_0\) common to all the embeddings.
\footnote{\href{https://github.com/Deicyde/riemann-hypothesis/blob/40e56bab/src/Development/AutomorphicArithmetic/AutomorphicVarieties/BuildingsResidueChars/HeckeCharParallelWeight.lean\#L39-L87}{Convergence of factorial powers of the finite unit idele},
\href{https://github.com/Deicyde/riemann-hypothesis/blob/40e56bab/src/Development/AutomorphicArithmetic/AutomorphicVarieties/BuildingsResidueChars/HeckeCharParallelWeight.lean\#L178-L287}{the minimum exponent and the choice of unit}.}

\subsection{Computing the jumps at conjugate embeddings}

Set \(b=a-1-s_0\). Equation~\eqref{eq:galois-inf-char} gives
\(I_{\sigma c}=\{a-1-b-z:z\in I_\sigma\}\). If
\(\iota:\overline{\mathbb Q}_\ell\to\mathbb C\) is the fixed
coefficient field isomorphism, define
\[
       H_\tau=\{(a-1)/2-z:z\in I_{\iota\tau}\}
\]
for each \(\tau:E\hookrightarrow\overline{\mathbb Q}_\ell\).
Each member is an integer by regular algebraicity. In the multiset
for \(\tau c\), the term belonging to \(a-1-b-z\) equals
\[
       (a-1)/2-(a-1-b-z)=b-\bigl((a-1)/2-z\bigr).
\]
This proves \(H_{\tau c}=\{b-h:h\in H_\tau\}\) with multiplicity.
For each local embedding in the scope of the de Rham comparison,
its jumps clause identifies \(H_\tau\) with the jumps at \(\tau\),
using exactly the displayed infinitesimal-character formula as its
hypothesis. Thus the same choice of \(b\) and \(H_\tau\) works at
each such place of the global representation.
\footnote{\href{https://github.com/Deicyde/riemann-hypothesis/blob/40e56bab/src/Development/AbelianVarietyTheory/Proofs/SatakeParameters/ExistsIsWeight.lean\#L138-L150}{The formula for the global weight},
\href{https://github.com/Deicyde/riemann-hypothesis/blob/40e56bab/src/Development/AutomorphicArithmetic/Foundations/GaloisRepresentations/IntegralStructures/InertCharpoly/GaloisLocalGlobalOfLocalGlobalCompatibility.lean\#L127-L171}{integral weights, conjugate weights and their occurrence in the local comparison}.}
